\chapter{Jenkins: CI server}
\label{ch:jenkins}

\section{Kiến trúc}

Jenkins là một ứng dụng web Java --- \emph{controller} --- lưu job,
credential, plugin và lịch sử build trong một thư mục,
\code{JENKINS\_HOME}, và thực thi build hoặc trên chính nó (qua các
\emph{executor}, các khe build song song của nó) hoặc trên các
\emph{agent} từ xa. Mẫu hình doanh nghiệp là một controller không
build gì cả, đẩy việc ra các agent dùng xong bỏ (thường là pod, qua
plugin Kubernetes).

\begin{decision}[build trên controller, Jenkins ngoài cluster]
Hai sự đảo ngược thực hành doanh nghiệp kiểu homelab, cả hai đều là
phép tính CPU. Agent pod sẽ tốn chi phí lập lịch và kéo image cho mỗi
lần build trên cùng hai nhân mà chính build cần; hai executor của
controller dành những chu kỳ đó cho Maven thay vì thế. Và Jenkins chạy
trong Docker trên máy chủ, không phải trong k3s: nếu bạn làm hỏng
cluster --- và trong lúc học, bạn sẽ làm hỏng --- thứ có thể dựng lại
cluster không được chết cùng nó. Cả hai đều đảo ngược gọn gàng ở quy
mô doanh nghiệp: Jenkins trong cluster với pod agent là nội dung của
Chương~27, và bản thân cuộc di dời chính là bài học.
\end{decision}

\section{Image của controller, có chú giải}

Thay vì bấm chuột cài công cụ vào một Jenkins nguyên bản, dự án nướng
toolchain vào một image --- \code{deploy/jenkins/Dockerfile}:

\begin{lstlisting}[language=dockerfile]
FROM jenkins/jenkins:lts-jdk21          (*@\co{1}@*)
USER root
RUN ... temurin-25-jdk ...              (*@\co{2}@*)
RUN ... apache-maven-3.9.9 ...          (*@\co{3}@*)
RUN wget -qO /usr/local/bin/kubectl ... (*@\co{4}@*)
ENV JDK25_HOME=/usr/lib/jvm/temurin-25-jdk-amd64  (*@\co{5}@*)
USER jenkins
RUN jenkins-plugin-cli --plugins \
      git workflow-aggregator ... credentials-binding  (*@\co{6}@*)
\end{lstlisting}

\begin{description}
  \item[\co{1}] Jenkins LTS chạy trên JDK\,21 --- phiên bản mà
  \emph{chính Jenkins} chạy trên đó.
  \item[\co{2}] Dự án nhắm tới Java\,25, nên một JDK thứ hai được cài
  cho việc \emph{build}. Lưu ý \code{JAVA\_HOME} \emph{không} bị đổi
  toàn cục \co{5} --- Jenkinsfile đặt nó cho từng pipeline. Runtime
  của controller và toolchain build là hai mối quan tâm độc lập;
  Dockerfile này giữ chúng tách bạch rõ ràng.
  \item[\co{3}] Maven ghim ở 3.9.9 --- gói của distro thường lạc hậu.
  \item[\co{4}] \code{kubectl}: stage deploy chỉ là binary này cộng
  với một kubeconfig.
  \item[\co{6}] Plugin ghim trong image, không bấm chọn trên UI: hỗ
  trợ pipeline, git, báo cáo JUnit, và \code{credentials-binding}
  (dùng bởi \code{withCredentials} ở Chương~23). Toàn bộ ý nghĩa:
  \emph{bản thân CI server giờ tái tạo được} --- hủy container, build
  lại từ Dockerfile, và chỉ \code{JENKINS\_HOME} (một volume có tên)
  mang trạng thái.
\end{description}

Một cờ lúc chạy trong \code{run.sh} xứng đáng có đoạn riêng:
\code{--network host}. Jib bên trong Jenkins push tới
\code{localhost:5000} --- và bên trong một container dùng mạng bridge,
\code{localhost} sẽ là \emph{container}. Mạng host làm Jenkins dùng
chung namespace mạng của máy chủ, nên \code{localhost:5000} của nó
chính là registry mà k3s kéo image từ đó (hợp đồng đặt tên của
Chương~13). Không phải đường tắt --- mà là yêu cầu của cách đặt tên
image.

\section{Credential}

Jenkins lưu secret trong một kho mã hóa, tham chiếu bằng ID; pipeline
gắn chúng vào biến môi trường chỉ trong phạm vi một khối
(\code{withCredentials}), và bộ che log xóa giá trị của chúng khỏi
đầu ra console. Dự án này lưu một credential: \code{kubeconfig-ts},
một credential dạng file chứa kubeconfig mà Jenkins dùng để triển khai
--- và đó là phần thú vị.

\section{Danh tính triển khai: một ServiceAccount, không phải admin}

Con đường lười biếng là đưa cho Jenkins kubeconfig admin của k3s ---
\code{cluster-admin}, được phép xóa mọi namespace. Thay vào đó,
\code{deploy/k8s/jenkins-rbac.yaml} dựng một danh tính tối thiểu:

\begin{lstlisting}[language=yaml]
kind: ServiceAccount          # the identity
metadata: { name: jenkins, namespace: ts }
---
kind: Role                    # namespaced permissions
metadata: { name: jenkins-deployer, namespace: ts }
rules:
  - apiGroups: ["apps"]
    resources: ["deployments", "statefulsets"]
    verbs: ["get", "list", "watch", "create", "patch", "update"] (*@\co{1}@*)
  - apiGroups: ["apps"]
    resources: ["replicasets"]
    verbs: ["get", "list", "watch"]                              (*@\co{2}@*)
  # + services, PVCs, secrets, configmaps, ingresses, pods/log ...
  - apiGroups: ["kafka.strimzi.io"]                              (*@\co{3}@*)
    resources: ["kafkas", "kafkanodepools", "kafkatopics"]
    verbs: ["get", "list", "create", "patch", "update"]
---
kind: RoleBinding             # identity <- permissions
\end{lstlisting}

\begin{description}
  \item[\co{1}] Các verb được suy ra từ chính các lệnh của pipeline:
  \code{kubectl apply} cần get+create+patch; \code{set image} cần
  get+patch. \emph{Không có delete ở bất cứ đâu}: CI chỉ tiến tới; con
  người mới xóa. Một pipeline bị chiếm quyền vì thế có thể làm hỏng
  các Deployment trong \code{ts} nhưng không thể quét sạch chúng, và
  không thể nhìn ra ngoài namespace --- Role, không phải ClusterRole,
  là quyết định về bán kính ảnh hưởng.
  \item[\co{2}] Phát hiện theo cách trung thực: \code{rollout status}
  theo dõi các ReplicaSet của Deployment, nên cần quyền đọc chúng ---
  loại quyền bạn tìm ra bằng cách chạy pipeline với Role đó và đọc
  lỗi \code{Forbidden} đầu tiên.
  \item[\co{3}] Thêm cùng cột mốc~6: pipeline được phép khai báo
  cluster và topic Kafka (custom resource theo namespace), nhưng cài
  operator và các CRD của nó --- phạm vi cluster --- vẫn là việc của
  admin (Chương~25). RBAC lớn thêm một rule cho mỗi API group mà
  pipeline chạm tới; file Role là nhật ký thay đổi của những gì CI đã
  được tin cậy.
\end{description}

Kubeconfig lưu trong Jenkins đơn giản ghép địa chỉ cluster với token
của ServiceAccount này thay vì chứng chỉ admin.

\begin{pitfall}[Forbidden: cannot create resource ``namespaces'']
Triệu chứng: stage deploy chết với lỗi Forbidden trên
\code{namespaces} dù file RBAC trông đã đầy đủ. Nguyên nhân: namespace
thuộc \emph{phạm vi cluster}, và một Role theo namespace không thể cấp
quyền đó --- chính xác là lý do \code{namespace.yaml} bị loại khỏi
kustomization (Chương~20) và namespace được admin tạo một lần. Gỡ lỗi
RBAC nói chung:
\code{kubectl auth can-i patch deployments -n ts --as
system:serviceaccount:ts:jenkins} trả lời ``việc này có được phép
không'' mà không cần chạy pipeline.
\end{pitfall}

\begin{handson}[chứng minh ranh giới tồn tại]
\begin{lstlisting}[language=bash]
kubectl auth can-i patch deployments -n ts \
  --as system:serviceaccount:ts:jenkins        # yes
kubectl auth can-i delete deployments -n ts \
  --as system:serviceaccount:ts:jenkins        # no
kubectl auth can-i list pods -n kube-system \
  --as system:serviceaccount:ts:jenkins        # no
\end{lstlisting}
\end{handson}

\begin{quiz}
\begin{enumerate}
  \item Những gì nằm trong \code{JENKINS\_HOME}, và vì sao nướng
  plugin vào image cộng với một volume cho home làm CI server trở
  thành thứ dùng xong bỏ được?
  \item Vì sao controller chạy hai JDK, và mỗi JDK được chọn ở đâu?
  \item Dựng lại danh sách verb RBAC cho một pipeline chỉ chạy
  \code{kubectl apply -k} và \code{rollout status} --- verb nào trong
  Role của dự án khi đó có thể bỏ đi?
  \item Giải thích \code{--network host} trên container Jenkins theo
  hợp đồng đặt tên image của Chương~13.
  \item ServiceAccount, Role, RoleBinding: cái nào là danh tính, cái
  nào là năng lực, cái nào là mối nối --- và vì sao
  Role-chứ-không-phải-ClusterRole là tiêu đề bảo mật?
\end{enumerate}
\end{quiz}
